(1) 发射圆环上每个微小电流元对轴线上的径向磁场分量互相抵消，轴向分量相加。Biot–Savart 积分给
\[
B_z(z)=\frac{\mu_0I_1a^2}{2(a^2+z^2)^{3/2}} .
\]
对 I₁>0 沿轴正方向为正；在 z=0 恰为 \(\mu_0I_1/(2a)\)。

(2) 若接收环半径 \(b\ll\sqrt{a^2+z^2}\)，其截面上的 Bz 可近似用中心值，磁通
\[
\Phi_{2\leftarrow1}\simeq\pi b^2 B_z
=\frac{\mu_0\pi a^2b^2}{2(a^2+z^2)^{3/2}}I_1.
\]
定义互感 \(\mathcal M=\Phi/I_1\)，故
\[
\mathcal M(z)\simeq\frac{\mu_0\pi a^2b^2}{2(a^2+z^2)^{3/2}}.
\]
遗漏的横向场非均匀修正通常是 \(b^2/(a^2+z^2)\) 的高阶量。

(3) 接收环闭合路径的感应电动势由总链接磁通的时间导数给出
\[
\mathcal E_2=-\frac{d}{dt}[\mathcal M(z(t))I_1(t)]
=-\mathcal M\dot I_1-I_1\mathcal M'(z)\dot z.
\]
其中
\(\mathcal M'(z)=-3\mu_0\pi a^2b^2z/[2(a^2+z^2)^{5/2}]\)。
第一项由磁源电流变化造成，第二项是运动穿过不均匀磁场造成的动生电动势；即使电流恒定，移动环仍可感应电动势。

(4) 接收圆环磁矩为 \(\boldsymbol\mu_2=I_2\pi b^2\hat z\)。在轴向变化的外磁场中固定两环电流时，偶极受力
\[
F_z=\mu_2\frac{dB_z}{dz}
=-\frac{3\mu_0I_1I_2\pi a^2b^2z}{2(a^2+z^2)^{5/2}}
=I_1I_2\mathcal M'(z).
\]
从势能 \(U=-\boldsymbol\mu_2\cdot\boldsymbol B_1=-I_1I_2\mathcal M\) 求 \(-U_z\) 得相同结果。远场 z≫a 时 \(F_z\sim-3\mu_0I_1I_2\pi a^2b^2/(2z^4)\)，同向电流环相互吸引。
